% SGA 6, Exposición VI: fragmento español alineado con la fuente.
% Alcance: sección 2, números 2.1--2.3.1, líneas 8907--8944 inclusive.
% Fuente francesa: sga6_fr_r1/src/fr.tex.
% SHA-256 de la fuente francesa completa:
% 0D2158B9FE3B466309A92165727A0B191422E05FBAA0DACD5806E8322C54BA1A.
\providecommand{\sisi}[2]{#2}

\section*{2. Cálculo del \(K'\) de un fibrado proyectivo: caso de los complejos perfectos}


Este parágrafo no se utilizará en lo que sigue del seminario. La notación \(K'(X)\) designará en este parágrafo el grupo de Grothendieck de los complejos perfectos sobre \(X\) (IV 2.2). Nos proponemos demostrar el teorema siguiente, análogo a 1.1:

\textbf{Teorema 2.1.} Sean \(S\) un esquema cuasicompacto, \(E\) un \(\mathcal O_S\)-módulo localmente libre de rango \(r+1\), y \(X=\mathbf P(E)\) el fibrado proyectivo asociado; sea \(\xi\) la clase en \(K'(X)\) del complejo igual a \(\mathcal O_X(1)\) en grado cero, y nulo en los demás grados. Entonces \(K'(X)\), dotado de su estructura natural de \(K'(S)\)-álgebra, es un \(K'(S)\)-módulo libre que admite como base \(1,\xi,\ldots,\xi^r\).

\textbf{2.2.} Fijemos las notaciones: \(B\) designará el álgebra simétrica de \(E\); denotaremos por \(\mathcal A\) la categoría de los \(\mathcal O_X\)-módulos, por \(\mathcal G\) la categoría de los \(B\)-módulos graduados cuyos morfismos son los morfismos graduados, por \(\mathcal G_0\) la subcategoría plena de \(\mathcal G\) cuyos objetos son los \(B\)-módulos localmente libres de tipo finito como \(B\)-módulos graduados, y por \(\mathcal G'\) la categoría de los \(\mathcal O_S\)-módulos. Las categorías \(\mathcal A\), \(\mathcal G\) y \(\mathcal G'\) son abelianas. Para todo \(n\), y todo complejo \(L^\bullet\) de \(D^+(\mathcal A)\), denotaremos por \(L^\bullet(n)\) el complejo \(L^\bullet\otimes_{\mathcal O_X}\mathcal O_X(n)\).

\textbf{2.3.} Necesitaremos aplicar el funtor \(\Proj_0\) a \(B\)-módulos graduados cualesquiera, y no solamente a \(B\)-módulos graduados cuasicoherentes como en la construcción de EGA II 3.2. Indicaremos, pues, rápidamente cómo se puede efectuar esta construcción en el caso general.

\textbf{2.3.1.} Sean primero \(S\) un esquema, \(A\) una \(\mathcal O_S\)-álgebra cuasicoherente, \(Y=\Spec(A)\), y \(f:Y\to S\) el morfismo estructural. Del isomorfismo
\[
  A\xrightarrow{\sim} f_*(\mathcal O_Y)
\]
se deduce un homomorfismo canónico:
\begin{equation}\tag{2.3.1.1}
  f^{-1}(A)\longrightarrow \mathcal O_Y,
\end{equation}
donde \(f^{-1}\) designa la imagen inversa en el sentido de los haces de conjuntos; en otros términos, se tiene un morfismo de espacios anillados:
\begin{equation}\tag{2.3.1.2}
  \bar f:(Y,\mathcal O_Y)\longrightarrow (S,A).
\end{equation}
Se deduce un funtor de «imagen inversa» \(\bar f^{\,*}\) de la categoría de los \(A\)-módulos en la categoría de los \(\mathcal O_Y\)-módulos, adjunto por la izquierda al funtor de «imagen directa» \(\bar f_*\) de la categoría de los \(\mathcal O_Y\)-módulos en la categoría de los \(f_*(\mathcal O_Y)\)-módulos. Se tiene, pues:
\begin{equation}\tag{2.3.1.3}
  \bar f^{\,*}(M)=f^{-1}(M)\otimes_{f^{-1}(A)}\mathcal O_Y
\end{equation}
para todo \(A\)-módulo \(M\).
Como el homomorfismo (2.3.1.1) es plano (pues en cada fibra corresponde a una localización), el funtor \(\bar f^{\,*}\) así definido es exacto (cf. [4], p. 164). Si \(M\) es un \(A\)-módulo cuasicoherente, esta construcción da el módulo \(\widetilde M\) definido en EGA II 1.4.3. En efecto, el isomorfismo
\[
  M\xrightarrow{\sim} f_*(\widetilde M)=\bar f_*(\widetilde M)
\]
da por adjunción un homomorfismo
\[
  \bar f^{\,*}(M)\longrightarrow \sisi{M}{\widetilde M}
\]
y se verifica inmediatamente en cada fibra que es un isomorfismo.

